\documentclass[11pt]{beamer}

\usetheme{Madrid}
\usecolortheme{default}

% Pacotes
\usepackage[utf8]{inputenc}
\usepackage[portuguese]{babel}
\usepackage{listings}
\usepackage{xcolor}
\usepackage{graphicx}
\usepackage{hyperref}

% Configuração de cores para código
\definecolor{codegreen}{rgb}{0,0.6,0}
\definecolor{codegray}{rgb}{0.5,0.5,0.5}
\definecolor{codepurple}{rgb}{0.58,0,0.82}
\definecolor{backcolour}{rgb}{0.95,0.95,0.92}

\lstset{
    backgroundcolor=\color{backcolour},
    commentstyle=\color{codegreen},
    keywordstyle=\color{codepurple},
    numberstyle=\tiny\color{codegray},
    stringstyle=\color{red},
    basicstyle=\ttfamily\tiny,
    breakatwhitespace=false,
    breaklines=true,
    captionpos=b,
    keepspaces=true,
    numbers=left,
    numbersep=5pt,
    showspaces=false,
    showstringspaces=false,
    showtabs=false,
    tabsize=2,
    language=bash,
    frame=single
}

% Título e Autor
\title{Introdução ao GitHub}
\subtitle{Para Usuários Windows}
\author{Seu Nome}
\date{\today}
\institute{Curso de Desenvolvimento}

\AtBeginSection[]
{
  \begin{frame}
    \frametitle{Sumário}
    \tableofcontents[currentsection]
  \end{frame}
}

\begin{document}

% Slide de Título
\begin{frame}
\titlepage
\end{frame}

% Slide de Sumário
\begin{frame}{Sumário}
\tableofcontents
\end{frame}

% ============================================================================
\section{O que é GitHub?}
% ============================================================================

\begin{frame}{O que é GitHub?}
\begin{itemize}
    \item Plataforma baseada em nuvem para hospedar código
    \item Usa \textbf{Git} como sistema de controle de versão
    \item Permite colaboração entre desenvolvedores
    \item Rastreia todas as mudanças no projeto
    \item Oferece backup automático
\end{itemize}

\begin{center}
    \large \textbf{Git} = Sistema local | \textbf{GitHub} = Hospedagem em nuvem
\end{center}
\end{frame}

\begin{frame}{Por que usar GitHub?}
\begin{columns}
\column{0.5\textwidth}
\textbf{Vantagens Técnicas:}
\begin{itemize}
    \small
    \item Controle de versão
    \item Histórico completo
    \item Reverter mudanças
    \item Branches independentes
\end{itemize}

\column{0.5\textwidth}
\textbf{Vantagens Profissionais:}
\begin{itemize}
    \small
    \item Portfólio online
    \item Colaboração em equipe
    \item Comunidade ativa
    \item Integração com ferramentas
\end{itemize}
\end{columns}
\end{frame}

% ============================================================================
\section{Conceitos Fundamentais}
% ============================================================================

\begin{frame}{Repositório (Repo)}
\begin{itemize}
    \item \textbf{Repositório} = Pasta de projeto com controle de versão
    \item Contém todos os arquivos + histórico de mudanças
    \item Pasta especial \texttt{.git} rastreia tudo
\end{itemize}

\vspace{0.5cm}

\textbf{Dois tipos:}
\begin{itemize}
    \item \textbf{Local}: No seu computador
    \item \textbf{Remoto}: No GitHub (na nuvem)
\end{itemize}

\vspace{0.5cm}

\begin{center}
    Seu PC (Local) $\longleftrightarrow$ GitHub (Remoto)
\end{center}
\end{frame}

\begin{frame}{Commit}
\textbf{Commit} = "Salvar" uma versão do seu projeto

\begin{itemize}
    \item É como fazer um backup com anotações
    \item Tem mensagem descritiva
    \item Registra QUEM fez, QUANDO e O QUÊ foi feito
    \item Cada commit tem ID único
\end{itemize}

\vspace{0.5cm}

\begin{center}
    \includegraphics[width=0.8\textwidth]{../../../common/images/commit-timeline.png}
\end{center}

\small (Se tiver imagem disponível)
\end{frame}

\begin{frame}{Exemplo de Commit}
\begin{lstlisting}
commit a1b2c3d4e5f6g7h8
Author: João Silva <joao@email.com>
Date:   Mon May 26 10:30:00 2026

    Adiciona função de login
    
    - Implementa autenticação de usuário
    - Adiciona validação de senha
    - Inclui testes unitários
\end{lstlisting}

\begin{itemize}
    \small
    \item \textbf{Hash}: Identificador único (a1b2c3d...)
    \item \textbf{Author}: Quem fez
    \item \textbf{Date}: Quando
    \item \textbf{Mensagem}: O quê e por quê
\end{itemize}
\end{frame}

\begin{frame}{Branch (Ramo)}
\textbf{Branch} = Linha de desenvolvimento independente

\begin{itemize}
    \item Trabalhe em features sem afetar o código principal
    \item \texttt{main} é o branch principal (produção)
    \item Crie branches como \texttt{feature/login}, \texttt{bugfix/error}
    \item Depois mescle (merge) de volta para main
\end{itemize}

\vspace{0.3cm}

\textbf{Exemplo:}
\begin{itemize}
    \small
    \item Branch \texttt{main}: Código em produção
    \item Branch \texttt{feature/2fa}: Desenvolvendo autenticação 2FA
    \item Quando pronto: Mescla para \texttt{main}
\end{itemize}
\end{frame}

\begin{frame}{Visualizando Branches}
\begin{lstlisting}
main:           A -> B -> C -> D -> E
                         |
feature/login:           -> F -> G
\end{lstlisting}

\begin{itemize}
    \item Trabalho isolado
    \item Sem conflitos
    \item Depois junta tudo: \texttt{git merge}
\end{itemize}
\end{frame}

\begin{frame}{Pull Request (PR)}
\textbf{Pull Request} = Solicitação formal para mesclar código

\begin{itemize}
    \item Você abre uma PR no GitHub
    \item Colegas \textbf{revisam} seu código
    \item Discussão e sugestões de melhorias
    \item Testes automatizados rodam
    \item Após aprovação: \textbf{Merge}
\end{itemize}

\vspace{0.5cm}

\textbf{Fluxo:}
\begin{enumerate}
    \small
    \item Faz mudanças no seu branch
    \item Push para GitHub
    \item Abre PR
    \item Time revisa
    \item Ajusta conforme feedback
    \item Aprova e mergeia
\end{enumerate}
\end{frame}

% ============================================================================
\section{Instalação no Windows}
% ============================================================================

\begin{frame}{Instalação: Passo 1}
\textbf{Baixar Git para Windows}

\begin{enumerate}
    \item Acesse \url{https://git-scm.com/}
    \item Clique em \textbf{Download for Windows}
    \item Escolha versão compatível (32 ou 64 bits)
    \item Execute o instalador (.exe)
\end{enumerate}

\vspace{0.5cm}

\begin{center}
    \textbf{Importante:} Fazer download oficial, não de sites desconhecidos
\end{center}
\end{frame}

\begin{frame}{Instalação: Passo 2}
\textbf{Instalador - Componentes Recomendados}

\begin{itemize}
    \item \textbf{Git Bash} ✓ (terminal para usar Git)
    \item \textbf{Git GUI} ✓ (interface visual)
    \item \textbf{Windows Explorer integration} ✓ (clicar direito)
    \item \textbf{Git Credential Manager} ✓
\end{itemize}

\vspace{0.3cm}

\textbf{Editor padrão:} Notepad ou VS Code

\vspace{0.3cm}

\textbf{Deixe outras opções no padrão}
\end{frame}

\begin{frame}{Instalação: Passo 3}
\textbf{Verificar Instalação}

Abra \textbf{Git Bash} ou \textbf{PowerShell}:

\begin{lstlisting}
git --version
\end{lstlisting}

Deve retornar:
\begin{lstlisting}
git version 2.45.0.windows.1
\end{lstlisting}

\vspace{0.5cm}

\textbf{✓ Pronto! Git instalado}
\end{frame}

\begin{frame}{Conta GitHub}
\textbf{Criar conta GitHub}

\begin{enumerate}
    \item Acesse \url{https://github.com/}
    \item Clique em \textbf{Sign up}
    \item Preencha dados:
    \begin{itemize}
        \item Email
        \item Senha (forte!)
        \item Username
    \end{itemize}
    \item Verifique seu email
\end{enumerate}

\vspace{0.3cm}

\textbf{Username dicas:}
\begin{itemize}
    \small
    \item Sem espaços
    \item Sem caracteres especiais (exceto - e \_)
    \item Profissional (é seu portfólio!)
\end{itemize}
\end{frame}

% ============================================================================
\section{Configuração Inicial}
% ============================================================================

\begin{frame}{Configurar Git}
\textbf{Primeira vez usando Git?}

Defina seu nome e email:

\begin{lstlisting}
git config --global user.name "João Silva"
git config --global user.email "joao@email.com"
\end{lstlisting}

\vspace{0.5cm}

\textbf{Verificar:}
\begin{lstlisting}
git config --list
\end{lstlisting}

\vspace{0.5cm}

\textbf{Nota:} Use a mesma email da conta GitHub!
\end{frame}

\begin{frame}{SSH Keys (Recomendado)}
\textbf{Por que SSH?} Não precisa digitar senha toda vez

\vspace{0.3cm}

\textbf{Gerar chave:}
\begin{lstlisting}
ssh-keygen -t ed25519 -C "seu@email.com"
\end{lstlisting}

Pressione Enter para tudo (usa padrões)

\vspace{0.5cm}

\textbf{Ver chave pública:}
\begin{lstlisting}
cat ~/.ssh/id_ed25519.pub
\end{lstlisting}

\vspace{0.3cm}

Copie a saída (começa com \texttt{ssh-ed25519})
\end{frame}

\begin{frame}{Adicionar Chave ao GitHub}
\begin{enumerate}
    \item Acesse \url{https://github.com/settings/keys}
    \item Clique \textbf{New SSH key}
    \item \textbf{Title}: "Meu Windows" (algo descritivo)
    \item \textbf{Key}: Cole a chave copiada
    \item Clique \textbf{Add SSH key}
\end{enumerate}

\vspace{0.5cm}

\textbf{Testar conexão:}
\begin{lstlisting}
ssh -T git@github.com
\end{lstlisting}

Deve retornar:
\begin{lstlisting}
Hi username! You've successfully authenticated...
\end{lstlisting}
\end{frame}

% ============================================================================
\section{Fluxo Básico}
% ============================================================================

\begin{frame}{Fluxo: Criar Repositório Local}
\textbf{Passo 1: Criar pasta e inicializar Git}

\begin{lstlisting}
mkdir meu-projeto
cd meu-projeto
git init
\end{lstlisting}

Cria pasta \texttt{.git} (oculta) que rastreia mudanças

\vspace{0.5cm}

\textbf{Passo 2: Criar alguns arquivos}

\begin{lstlisting}
echo "# Meu Projeto" > README.md
\end{lstlisting}

Ou use VS Code:
\begin{lstlisting}
code .
\end{lstlisting}
\end{frame}

\begin{frame}{Fluxo: Verificar Status}
\textbf{Quais arquivos foram modificados?}

\begin{lstlisting}
git status
\end{lstlisting}

\vspace{0.3cm}

Retorna:
\begin{lstlisting}
On branch main
Untracked files:
  (use "git add <file>..." to include in what will be committed)
	README.md

nothing added to commit but untracked files present
\end{lstlisting}

\vspace{0.3cm}

\textbf{Status possíveis:}
\begin{itemize}
    \small
    \item \textcolor{red}{Untracked}: Arquivo novo, Git não conhece
    \item \textcolor{red}{Modified}: Arquivo existente foi mudado
    \item \textcolor{green}{Staged}: Pronto para commit
\end{itemize}
\end{frame}

\begin{frame}{Fluxo: Staging Area}
\textbf{Adicionar arquivos à "staging area"}

Antes de commitar, preparamos os arquivos:

\begin{lstlisting}
git add README.md
\end{lstlisting}

Ou adicionar tudo:
\begin{lstlisting}
git add .
\end{lstlisting}

\vspace{0.5cm}

\textbf{Verificar novamente:}
\begin{lstlisting}
git status
\end{lstlisting}

Agora mostra em verde (pronto para commit)
\end{frame}

\begin{frame}{Fluxo: Fazer Commit}
\textbf{Criar snapshot do projeto}

\begin{lstlisting}
git commit -m "Adiciona README inicial"
\end{lstlisting}

\vspace{0.5cm}

\textbf{Boa mensagem de commit:}
\begin{itemize}
    \item Começa com verbo no imperativo (Adiciona, Fixa, Atualiza)
    \item Breve e descritiva
    \item Primeira letra maiúscula
    \item Sem ponto final
\end{itemize}

\vspace{0.5cm}

\textbf{Exemplos:}
\begin{itemize}
    \small
    \item ✓ "Adiciona validação de email"
    \item ✗ "email validation"
    \item ✗ "fix stuff"
\end{itemize}
\end{frame}

\begin{frame}{Fluxo: Ver Histórico}
\textbf{Ver todos os commits}

\begin{lstlisting}
git log
\end{lstlisting}

Mostra histórico completo (pressione Q para sair)

\vspace{0.3cm}

\textbf{Versão resumida:}
\begin{lstlisting}
git log --oneline
\end{lstlisting}

\vspace{0.3cm}

Retorna algo como:
\begin{lstlisting}
a1b2c3d Adiciona README inicial
f4e5d6c Adiciona .gitignore
\end{lstlisting}
\end{frame}

% ============================================================================
\section{Repositório Remoto}
% ============================================================================

\begin{frame}{Criar Repo no GitHub}
\begin{enumerate}
    \item Acesse \url{https://github.com/new}
    \item \textbf{Repository name}: \texttt{meu-projeto}
    \item \textbf{Description}: (opcional)
    \item Escolha \textbf{Public} ou \textbf{Private}
    \item \textbf{NÃO} marque "Initialize with README"
    \item Clique \textbf{Create repository}
\end{enumerate}

\vspace{0.3cm}

GitHub vai mostrar comandos para conectar seu repo local
\end{frame}

\begin{frame}{Conectar Local ao Remoto}
\textbf{Linking seu repositório local com GitHub}

\begin{lstlisting}
git remote add origin git@github.com:seu-usuario/meu-projeto.git
\end{lstlisting}

\textbf{Renomear branch (Git moderno):}
\begin{lstlisting}
git branch -M main
\end{lstlisting}

\textbf{Enviar commits para GitHub:}
\begin{lstlisting}
git push -u origin main
\end{lstlisting}

\vspace{0.3cm}

Agora seu código está no GitHub! 🎉
\end{frame}

\begin{frame}{Depois: Atualizar}
\textbf{Após fazer mais commits locais}

Simplesmente:
\begin{lstlisting}
git push
\end{lstlisting}

Envia novos commits para GitHub

\vspace{0.5cm}

\textbf{Baixar mudanças do GitHub:}
\begin{lstlisting}
git pull
\end{lstlisting}

Traz commits feitos por outros em seu PC
\end{frame}

% ============================================================================
\section{Branches e PRs}
% ============================================================================

\begin{frame}{Criar um Branch}
\textbf{Trabalhar em uma nova feature}

\begin{lstlisting}
git checkout -b feature/login
\end{lstlisting}

Cria novo branch e muda para ele

\vspace{0.5cm}

\textbf{Verificar branches:}
\begin{lstlisting}
git branch
\end{lstlisting}

Mostra: \texttt{* feature/login} (asterisco = atual)

\vspace{0.5cm}

\textbf{Trocar de branch:}
\begin{lstlisting}
git checkout main
\end{lstlisting}
\end{frame}

\begin{frame}{Commits no Branch}
\textbf{Trabalhe normalmente}

\begin{lstlisting}
# Editar arquivos
# (abra no VS Code e modifique)

# Adicionar e commitar
git add .
git commit -m "Implementa tela de login"
git commit -m "Adiciona validação de email"
\end{lstlisting}

\vspace{0.5cm}

\textbf{Branch agora tem commits novos!}

main fica inalterado, seguro
\end{frame}

\begin{frame}{Push do Branch}
\textbf{Enviar branch para GitHub}

\begin{lstlisting}
git push origin feature/login
\end{lstlisting}

GitHub cria branch remoto automaticamente

\vspace{0.5cm}

\textbf{Ver todos os branches:}
\begin{lstlisting}
git branch -a
\end{lstlisting}

Mostra locais e remotos
\end{frame}

\begin{frame}{Abrir Pull Request}
\begin{enumerate}
    \item Acesse seu repositório no GitHub
    \item Clique em \textbf{Pull requests}
    \item Clique \textbf{New pull request}
    \item Selecione:
    \begin{itemize}
        \item \textbf{Base}: \texttt{main}
        \item \textbf{Compare}: \texttt{feature/login}
    \end{itemize}
    \item Clique \textbf{Create pull request}
\end{enumerate}
\end{frame}

\begin{frame}{Descrever a PR}
\textbf{Preenchendo a descrição}

\begin{itemize}
    \item \textbf{Title}: Resumo da mudança
    \item \textbf{Description}: Detalhes (use Markdown!)
\end{itemize}

\vspace{0.3cm}

\textbf{Template útil:}

\begin{itemize}
    \small
    \item O que faz essa PR?
    \item Por que é necessária?
    \item Como testar?
    \item Screenshots (se visual)
    \item Issues relacionadas
\end{itemize}
\end{frame}

\begin{frame}{Revisar e Mesclar}
\textbf{Como revisor:}

\begin{itemize}
    \item Leia descrição
    \item Clique em \textbf{Files changed}
    \item Deixe comentários
    \item Aprove ou solicite mudanças
\end{itemize}

\vspace{0.5cm}

\textbf{Como autor (respondendo feedback):}

\begin{lstlisting}
git add .
git commit -m "Implementa feedback da revisão"
git push origin feature/login
\end{lstlisting}

PR se atualiza automaticamente

\vspace{0.3cm}

\textbf{Após aprovação:} Clique \textbf{Merge pull request}
\end{frame}

\begin{frame}{Depois de Mesclar}
\textbf{Limpeza}

\begin{lstlisting}
# Voltar para main
git checkout main

# Atualizar main local
git pull origin main

# Deletar branch local
git branch -d feature/login

# Deletar branch remoto (GitHub oferece botão)
git push origin --delete feature/login
\end{lstlisting}

\vspace{0.3cm}

\textbf{Pronto! Feature integrada em produção}
\end{frame}

% ============================================================================
\section{Comandos Essenciais}
% ============================================================================

\begin{frame}{Comandos: Consulta}
\begin{tabular}{|l|l|}
\hline
\textbf{Comando} & \textbf{O que faz} \\
\hline
\texttt{git status} & Status atual dos arquivos \\
\texttt{git log} & Histórico de commits \\
\texttt{git log --oneline} & Histórico resumido \\
\texttt{git diff} & Mudanças não commitadas \\
\texttt{git branch} & Lista branches locais \\
\texttt{git branch -a} & Todos os branches \\
\hline
\end{tabular}
\end{frame}

\begin{frame}{Comandos: Modificação}
\begin{tabular}{|l|l|}
\hline
\textbf{Comando} & \textbf{O que faz} \\
\hline
\texttt{git add <file>} & Prepara arquivo \\
\texttt{git add .} & Prepara tudo \\
\texttt{git commit -m "msg"} & Cria commit \\
\texttt{git reset HEAD <file>} & Remove de staging \\
\texttt{git checkout <file>} & Descarta mudanças \\
\hline
\end{tabular}
\end{frame}

\begin{frame}{Comandos: Remoto}
\begin{tabular}{|l|l|}
\hline
\textbf{Comando} & \textbf{O que faz} \\
\hline
\texttt{git push} & Envia para GitHub \\
\texttt{git pull} & Baixa do GitHub \\
\texttt{git fetch} & Baixa sem mesclar \\
\texttt{git clone <url>} & Copia repositório \\
\texttt{git remote -v} & Ver remotes \\
\hline
\end{tabular}
\end{frame}

% ============================================================================
\section{Boas Práticas}
% ============================================================================

\begin{frame}{Mensagens de Commit}
\textbf{Bom exemplo:}
\begin{lstlisting}
feat: Adiciona autenticação 2FA

- Implementa módulo de 2FA
- Integra com Twilio
- Adiciona testes
\end{lstlisting}

\vspace{0.5cm}

\textbf{Ruim:}
\begin{lstlisting}
fix
asdfgh
mudanças
\end{lstlisting}

\vspace{0.5cm}

\textbf{Prefixos úteis:}
\begin{itemize}
    \small
    \item \texttt{feat:} Nova funcionalidade
    \item \texttt{fix:} Correção
    \item \texttt{docs:} Documentação
    \item \texttt{refactor:} Reorganização
    \item \texttt{test:} Testes
\end{itemize}
\end{frame}

\begin{frame}{.gitignore}
\textbf{Arquivo que lista o que GIT deve IGNORAR}

\begin{lstlisting}
# .gitignore
node_modules/
__pycache__/
.env
*.log
.DS_Store
.vscode/
\end{lstlisting}

\begin{itemize}
    \item Não commita arquivos temporários
    \item Não commita senhas/chaves
    \item Não commita dependências
\end{itemize}
\end{frame}

\begin{frame}{Estrutura de Branches}
\textbf{Padrão comum em times:}

\begin{itemize}
    \item \texttt{main}: Produção (sempre estável)
    \item \texttt{develop}: Integração
    \item \texttt{feature/*}: Novas features
    \item \texttt{bugfix/*}: Correções
    \item \texttt{hotfix/*}: Urgente em produção
\end{itemize}

\vspace{0.5cm}

\textbf{Fluxo:}
\begin{enumerate}
    \small
    \item Feature branchs saem de \texttt{develop}
    \item PRs vão para \texttt{develop}
    \item Periodicamente: \texttt{develop} → \texttt{main}
\end{enumerate}
\end{frame}

% ============================================================================
\section{Resolução de Problemas}
% ============================================================================

\begin{frame}{Problema: SSH Permission Denied}
\textbf{Erro: "Permission denied (publickey)"}

\begin{lstlisting}
# Testar SSH
ssh -T git@github.com

# Adicionar chave ao ssh-agent
eval "$(ssh-agent -s)"
ssh-add ~/.ssh/id_ed25519

# Testar novamente
ssh -T git@github.com
\end{lstlisting}

\vspace{0.3cm}

Se ainda falhar: Reconfigurar SSH (ver README)
\end{frame}

\begin{frame}{Problema: Mensagem de Commit Errada}
\textbf{Corrigir último commit}

\begin{lstlisting}
git commit --amend -m "Nova mensagem"
\end{lstlisting}

\vspace{0.5cm}

\textbf{Incluir arquivo esquecido:}

\begin{lstlisting}
git add arquivo-esquecido.txt
git commit --amend --no-edit
\end{lstlisting}

\vspace{0.3cm}

\textbf{Cuidado:} Só faça antes de fazer push!
\end{frame}

\begin{frame}{Problema: Merge Conflict}
\textbf{Quando dois branches modificam mesma linha}

Git vai avisar:
\begin{lstlisting}
CONFLICT (content merge): Merge conflict in arquivo.txt
\end{lstlisting}

\vspace{0.3cm}

\textbf{Ver conflito:}
\begin{lstlisting}
<<<<<<< HEAD
seu código
=======
código outro
>>>>>>> branch
\end{lstlisting}

Escolha qual código manter, salve, e:
\begin{lstlisting}
git add arquivo.txt
git commit -m "Resolve merge conflict"
\end{lstlisting}
\end{frame}

\begin{frame}{Problema: Não é um Git Repo}
\textbf{Erro: "fatal: not a git repository"}

Você está fora do repositório

\vspace{0.3cm}

\textbf{Solução:}
\begin{lstlisting}
# Navegue para pasta correta
cd caminho/do/repositorio

# Ou inicie novo repo
git init
\end{lstlisting}
\end{frame}

% ============================================================================
\section{Próximos Passos}
% ============================================================================

\begin{frame}{Próximos Passos}
\begin{enumerate}
    \item \textbf{Configure Git} no seu Windows
    \item \textbf{Crie uma conta} GitHub
    \item \textbf{Clone um projeto} existente
    \item \textbf{Crie seu primeiro repo} pessoal
    \item \textbf{Pratique} commits, branches, PRs
    \item \textbf{Contribua} em projetos open-source
\end{enumerate}

\vspace{0.5cm}

Confira o README.md para mais detalhes!
\end{frame}

\begin{frame}{Recursos Úteis}
\begin{itemize}
    \item \url{https://git-scm.com/doc} - Documentação oficial Git
    \item \url{https://docs.github.com} - Documentação GitHub
    \item \url{https://learngitbranching.js.org} - Aprender Git interativamente
    \item \url{https://skills.github.com} - Tutoriais oficiais GitHub
\end{itemize}

\vspace{0.5cm}

\begin{center}
    \Large \textbf{Dúvidas?}
    
    \vspace{0.3cm}
    
    Abra uma issue no repositório!
\end{center}
\end{frame}

\begin{frame}
\begin{center}
    \Huge \textbf{Obrigado!}
    
    \vspace{0.5cm}
    
    \Large Bom aprendizado com Git e GitHub! 🚀
\end{center}
\end{frame}

\end{document}
